\section{Setting and Threat Model}
\label{sec:setting}

Several agents, typically different LLMs, answer the same question, each with an answer, a stated confidence, and citations to the passages they used. An \emph{arbitration step} decides the combined answer. We want that decision to have three properties an auditor can check. \emph{Replay}: from the recorded inputs, anyone can recompute the same decision bit for bit. \emph{Provenance}: the decision is one of the submitted answers, and every input it read traces to an agent, a captured passage, or a recorded checker label. \emph{An explicit rule}: why one answer beat another is a stated function of recorded inputs, so a changed rule can be re-applied to old decisions and compared.

Before selection, the answers, cited passage captures, checker labels, and policy inputs are recorded in a sealed artifact set. The reducer then reads only those recorded inputs, rather than fetching passages or asking a model again. Content hashes let a verifier detect edits against an independently held commitment, subject to hash-collision assumptions; they do not establish that a passage or label is true. \suppref{S9} specifies capture, sealing, and the replay guarantee.

These properties suit applications that need to inspect and revise decisions after their first use. An auditor can inspect the passage that licensed an overturn, remove an agent, or compare selectors using the same sealed inputs. New answers, passages, proposition renderings, or checker identities can require new labels. The experiments therefore compare both first-decision quality and the cost of later permitted re-decisions (\cref{sec:overturn-results,sec:audit-results}); reproducibility alone does not imply better accuracy.

The primary threat is the \emph{honest-but-fallible} agent, which hallucinates, holds stale knowledge, or is miscalibrated. The case verify-then-vote targets is \emph{correlated error}, where most agents assert the same wrong answer. The trusted computing base is the reducer, the capture tooling, and the checker. A checker can err, and a retrieved source can itself assert a wrong answer. The property in \cref{sec:gate} identifies these verified wrong answers as necessary for losing a correct default; it supplies no numerical guarantee on their frequency. An agent may also cite a genuine source that says something other than what it claims. That leaves every hash intact, and it is exactly what the checker is intended to catch.
